% Energy estimates and a priori bounds
% Converted from notes/pde/energy-estimates.html. Edit freely: this file is now the master copy.

\section{Energy estimates and a priori bounds}\label{p:energy-estimates}

\prereqs{Hölder, Young, and L\textsuperscript{\textit{p}} interpolation~(\S\ref{p:elementary-inequalities}), Gronwall's inequality and continuity arguments~(\S\ref{p:gronwall})}

\subsection*{The idea}

An a priori estimate is a bound that every smooth solution must satisfy, proved before knowing that solutions exist. The standard way to find one: multiply the equation by a well-chosen quantity, integrate, integrate by parts, and use the structure of the equation so that the dangerous terms cancel or have a sign. The choice of multiplier is the idea; the rest is the inequality toolkit~(\S\ref{p:elementary-inequalities}).

\subsection*{Statement}

\begin{thmbox}{The method (not a theorem)}
\begin{enumerate}
\item Multiply each equation by a multiplier (often the unknown itself, its time derivative, or $\Delta$ of it) and integrate over $\Omega$.
\item Integrate by parts, using the boundary conditions, to produce $\frac{d}{dt}$(energy) plus nonnegative dissipation.
\item Arrange the multipliers so that coupling terms cancel between equations.
\item Bound what remains by the energy and dissipation (Hölder, Sobolev, Young), absorb the dissipation parts, and close with Gronwall~(\S\ref{p:gronwall}).
\item For higher regularity, repeat after differentiating the equation in $t$ or $x$.
\end{enumerate}
\end{thmbox}

\begin{exbox}{Model example (heat equation)}
For $u_t-\Delta u=0$ with $u=0$ on $\partial\Omega$, multiplying by $u$ gives $\frac{d}{dt}\frac12\norm{u}_{L^2}^2+\norm{\nabla u}_{L^2}^2=0$, and multiplying by $u_t$ gives $\norm{u_t}_{L^2}^2+\frac{d}{dt}\frac12\norm{\nabla u}_{L^2}^2=0$.
\end{exbox}

\subsection*{Why it is plausible}

Physical systems usually dissipate energy, and the mathematical energy is the physical one. Kinetic energy, internal (pressure) energy and potential energy together cannot grow; each piece then bounds a norm of the unknowns. When the energy contains a term with no sign, the first job is to show it is still bounded below, typically by interpolation against a conserved quantity.

\subsection*{Exercises}

\begin{exercise}
  Derive the two identities for the heat equation in the model example, with every integration by parts and boundary term. Conclude that $u_0 \in H^1_0$ gives $u \in L^\infty(0,T;H^1_0)$ and $u_t \in L^2(0,T;L^2)$.
  \begin{hint}
    $\int u_t u = \frac12\frac{d}{dt}\norm{u}^2$ and $-\int\Delta u\,u = \norm{\nabla u}^2$ (boundary term vanishes since $u = 0$). Similarly $-\int\Delta u\,u_t = \frac12\frac{d}{dt}\norm{\nabla u}^2$ because $u_t = 0$ on $\partial\Omega$. Integrate in time.
  \end{hint}
  \begin{solution}
    Multiplying by $u$: $\int u_tu = \frac12\frac{d}{dt}\norm{u}^2$ and $-\int\Delta u\,u = \int|\nabla u|^2 - \int_{\partial\Omega}u\,\partial_\nu u = \norm{\nabla u}^2$ since $u = 0$ on $\partial\Omega$. So $\frac12\frac{d}{dt}\norm{u}^2 + \norm{\nabla u}^2 = 0$.

    Multiplying by $u_t$: $-\int\Delta u\,u_t = \int\nabla u\cdot\nabla u_t - \int_{\partial\Omega}u_t\,\partial_\nu u = \frac12\frac{d}{dt}\norm{\nabla u}^2$, because $u_t = 0$ on $\partial\Omega$ ($u$ vanishes there for all $t$). So $\norm{u_t}^2 + \frac12\frac{d}{dt}\norm{\nabla u}^2 = 0$.

    Integrating: $\norm{u(t)}^2 \le \norm{u_0}^2$ and $\norm{\nabla u(t)}^2 + 2\int_0^t\norm{u_t}^2 = \norm{\nabla u_0}^2$. So $u \in L^\infty(0,T;H^1_0)$ and $u_t \in L^2(0,T;L^2)$ when $u_0 \in H^1_0$.
  \end{solution}
\end{exercise}

\begin{exercise}
  A nonlinear term with a sign. For $u_t - \Delta u = u - u^3$ with $\partial_\nu u = 0$, multiply by $u$ and show $\norm{u(t)}_{L^2}^2 \le \norm{u_0}_{L^2}^2 + \frac{|\Omega|}{2}t$.
  \begin{hint}
    $\frac12\frac{d}{dt}\norm{u}^2 + \norm{\nabla u}^2 = \int(u^2 - u^4)$, and $s^2 - s^4 \le \frac14$ for all real $s$.
  \end{hint}
  \begin{solution}
    Multiplying by $u$ and using $\partial_\nu u = 0$: $\frac12\frac{d}{dt}\norm{u}^2 + \norm{\nabla u}^2 = \int_\Omega(u^2 - u^4)$. Since $s^2 - s^4 = \frac14 - (s^2 - \frac12)^2 \le \frac14$, the right side is at most $\frac{|\Omega|}{4}$. Dropping $\norm{\nabla u}^2 \ge 0$, $\frac{d}{dt}\norm{u}^2 \le \frac{|\Omega|}{2}$, and integrating, $\norm{u(t)}^2 \le \norm{u_0}^2 + \frac{|\Omega|}{2}t$. The cubic term was not estimated at all: its sign did the work.
  \end{solution}
\end{exercise}

\begin{exercise}
  Choosing the multiplier. For the wave equation $u_{tt} = u_{xx}$ on $(0,1)$ with $u = 0$ at both ends, find a multiplier that shows $\int_0^1 (u_t^2 + u_x^2)\,dx$ is constant in time.
  \begin{hint}
    Multiply by $u_t$: $u_{tt}u_t = \frac12(u_t^2)_t$ and $\int u_{xx}u_t = -\int u_x u_{xt} = -\frac12\frac{d}{dt}\int u_x^2$, since $u_t = 0$ at the ends.
  \end{hint}
  \begin{solution}
    Multiply by $u_t$ and integrate over $(0,1)$. First, $\int u_{tt}u_t = \frac12\frac{d}{dt}\int u_t^2$. Second, $\int u_{xx}u_t = \big[u_xu_t\big]_0^1 - \int u_xu_{xt} = -\frac12\frac{d}{dt}\int u_x^2$, since $u_t = 0$ at $x = 0, 1$. So $\frac{d}{dt}\int_0^1(u_t^2 + u_x^2)\,dx = 0$.
  \end{solution}
\end{exercise}

\subsection*{Where to find a proof}

\begin{itemize}
\item L. C. Evans, \emph{Partial Differential Equations}, 2nd ed., AMS, 2010, §2.3.4 and §7.1.2 (energy methods for the heat equation and for linear parabolic equations).
\end{itemize}

\neededfor{The Faedo--Galerkin method~(\S\ref{p:galerkin}), Uniqueness by difference estimates~(\S\ref{p:uniqueness-energy}), Regularized approximation schemes~(\S\ref{p:regularized-approximation})}
